% =====================================================================
% sec-intro.tex — 第 1 节 引言
% 1.1 研究背景与科学问题界定
% 1.2 国内外研究现状综述（双脉络文献演进与既有研究统计缺口）
% 1.3 研究目标与主要工作
% 1.4 论文结构安排
% =====================================================================
\section{引言}\label{sec:intro}

\subsection{研究背景与科学问题界定}

随着深度神经网络在计算机视觉、自然语言处理等关键领域的突破性进展，深度学习模型已被广泛部署于医疗辅助诊断、自动驾驶避障决策、金融高频风控及司法量刑辅助等一系列安全攸关（Safety-Critical）的高风险现实场景中。在这些高风险决策体系中，算法系统所面临的核心诉求不仅局限于提供具有高判别精度的离散预测类别 $\hat Y \in \{1, \dots, C\}$，更为关键的是要求模型输出的后验预测概率向量 $\hat{\boldsymbol{p}} = (\hat p_1, \dots, \hat p_C)^\top$ 具备高度可靠的不确定性度量能力。

从数理统计的严格视角来看，预测概率的可靠性体现为统计\textbf{条件概率校准性}（Conditional Calibration）。形式化地，设输入特征空间为 $\mathcal{X} \subset \mathbb{R}^d$，真实类别标签空间为 $\mathcal{Y} = \{1, \dots, C\}$，其联合概率测度记为 $P(X, Y)$。对于任意给定的分类模型，记其离散预测标签为 $\hat Y(X) = \operatorname{argmax}_{c \in \{1,\dots,C\}} \hat p_c(X)$，所对应的最大预测置信度为 $\hat P(X) = \max_{c \in \{1,\dots,C\}} \hat p_c(X) \in [1/C, 1]$。若对于所有可能的置信度水平 $p \in [0, 1]$，真实类别的后验发生概率与模型输出的预测置信度严格相等，即满足：
\begin{equation}\label{eq:strong_calib_intro}
\mathbb{P}\big(Y = \hat Y(X) \;\big|\; \hat P(X) = p\big) = p, \quad \forall p \in [0, 1],
\end{equation}
则称该模型在置信度意义下满足严格条件概率校准。若式~\eqref{eq:strong_calib_intro}成立，则当模型以 $95\%$ 的置信度预测某批样本为特定类别时，该批样本中恰有 $95\%$ 为分类正确，下游决策引擎即可据此设定最优拒绝推断阈值或风险对冲策略。

然而，现代深度学习的发展却伴随着这一概率度量能力的系统性退化。自残差网络（ResNet）、批归一化（Batch Normalization）及大容量卷积架构普及以来，尽管模型在各类基准测试上的 Top-1 错误率被不断刷新，但输出概率却呈现出普遍而严重的「过度自信」（Overconfidence）倾向。在经验交叉熵损失函数的持续驱动下，梯度下降优化算法在模型分类准确率饱和后，仍会持续驱使输出对数几率（Logits）的模长不断增大，进而导致 Softmax 函数将预测概率极化推向概率单纯形的极角（多数样本置信度集中于 $0.95$ 以上）。尤为严峻的是，当深层网络做出错误决策时，其平均输出置信度仍往往高达 $80\%$ 甚至更高，形成严重的「虚假高确信」现象。

从统计推断的理论视角来看，过度自信在数学上表现为深度模型参数化输出分布与经验真实条件分布的系统性偏差。针对这一科学问题，学术界提出了以后处理方式缩放 logits 向量模长的 Temperature Scaling 方法。然而，该方法是否在数理统计意义上具备严密的参数估计基础？其温度估计量是否具有良好的渐近性质？如何构建规范的假设检验体系以判定网络是否显著偏离校准状态？网络容量深度对校准程度与统计推断结构有何内在影响？这些关键问题构成了本项研究探索的核心科学议题。

\subsection{国内外研究现状综述}

围绕分类模型的概率校准与不确定性评估，统计学界与机器学习界历经数十年的交叉演进，形成了两条层次鲜明的理论与方法脉络。

\textbf{脉络一：统计适当评分规则与经典概率校准理论。}
概率预测评估的理论基石源于气象统计与贝叶斯推断中关于「适当评分规则」（Proper Scoring Rules）的经典研究。Brier~\cite{brier1950verification} 首先提出了基于均方误差度量的 Brier Score，开启了概率预测准确性定量检验的先河。此后，Gneiting 与 Raftery~\cite{gneiting2007strictly} 在《美国统计学会会刊》（\textit{JASA}）上发表系统性综述，确立了严格适当评分规则（Strictly Proper Scoring Rules）的公理化数学框架，在测度论层面严格证明了：对数似然损失（Negative Log-Likelihood, NLL）与 Brier Score 的期望泛函取得全局极小值，当且仅当模型的预测条件分布与真实数据生成分布完全一致。该定理从数理统计底层证明了优化似然函数等价于追求最优概率校准。在传统统计学习中，针对支持向量机（SVM）等非概率边际分类器，Platt~\cite{platt1999probabilistic} 提出了基于二项似然的单变量双参数 Logistic 校准模型（即经典的 Platt Scaling）。随后，Zadrozny 与 Elkan~\cite{zadrozny2002transforming} 进一步引入了基于保序回归（Isotonic Regression）的非参数单调校准算法及直方图分箱映射，确立了经典机器学习后处理校准的基本范式。

\textbf{脉络二：现代深度神经网络的校准偏差与缩放校准方法。}
伴随着现代深度神经网络在参数深度与宽度上的大幅扩张，经典校准假设遭遇了显著挑战。Guo 等~\cite{guo2017calibration} 在国际机器学习大会（\textit{ICML}）上系统报告了现代深层网络的校准偏差现象：与早期的浅层多层感知机（MLP）或支持向量机相比，现代深层网络（如 ResNet、DenseNet）虽然判别精度大幅提高，但其条件校准误差却发生了系统性恶化，呈现出显著的过度自信；为此，他们提出了引入标量温度 $T$ 的 Temperature Scaling 后处理方案，在经验评估中展示了良好的分箱误差压缩效果。随后，Vaicenavicius 等~\cite{vaicenavicius2019evaluating} 从多变量分布检验视角探讨了连续校准误差度量的统计假设检验与核方法扩展；Minderer 等~\cite{minderer2021revisiting} 则在 \textit{NeurIPS} 上系统考察了视觉 Transformer（ViT）与卷积模型在跨领域分布偏移下的校准漂移现象，指出模型架构与训练目标对不确定性度量具有决定性约束。

\textbf{既有研究的统计推断缺口。}
综观既有文献，尽管以 Guo 等为代表的现代机器学习研究广泛应用了 Temperature Scaling，但绝大多数工作仍将其视作经验性的后处理超参数调优，尚缺乏系统性的数理统计学理论分析。具体表现在：
\begin{enumerate}
\item 鲜有文献从指数族极大似然估计角度，对温度参数（或其倒数逆温度 $\beta=1/T$）的对数似然严格凹性、解析 Score 方程与 Fisher 信息矩阵给出完备的数学证明，缺乏对估计解全局唯一存在性的理论保证；
\item 缺少基于大样本渐近理论（Asymptotic Normality）与 Delta 方法对温度参数估计量 $\hat T$ 渐近方差的解析推导，缺乏与非参数 Bootstrap 的有限样本区间覆盖对比；
\item 缺乏针对「网络是否需要校准」这一核心命题（即原假设 $H_0: T=1$）的形式化假设检验体系（Wald 检验、Score 检验与似然比检验 LRT），缺少对有限样本下三大检验统计量因似然曲率变化导致数值排序差异的数理阐释；
\item 缺乏结合网络深度对比验证估计量 $\sqrt{n}$ 渐近相合速率与有限总体修正（FPC）机制的精细实证支撑。
\end{enumerate}
针对上述学术空白，本文旨在将 Temperature Scaling 全面纳入经典数理统计理论体系，通过严谨的公式推导、渐近分析与无泄漏对照实验，构建一套完整的极大似然推断与假设检验分析框架。

\subsection{研究目标与主要工作}

本文以现代残差网络架构在标准 CIFAR-10 图像分类任务上的后验概率为实证载体，以浅层 ResNet-20（20 层，0.27M 参数）与深层 ResNet-56（56 层，0.85M 参数）为对照实体，将经典的数理统计推断与假设检验理论系统应用于已有的概率校准方法，并开展无泄漏实证对比分析。本文的主要工作概括如下：

\begin{enumerate}
\item \textbf{已有校准模型的指数族似然整理与 Lean 4 形式化核验}：针对 Guo 等~\cite{guo2017calibration} 提出的单参数 Temperature Scaling 方法，将其整理表达为多分类 Softmax 指数族条件似然模型，显式写出其一阶 Score 函数与二阶观测 Fisher 信息量（即类别 logit 的条件方差之和），阐明似然函数的严格凹性与 MLE 解的唯一存在性，采用标准的阻尼 Newton--Raphson 迭代法进行数值求解，并将上述核心推导在 Lean 4（Mathlib）环境下完成了机器形式化证明。
\item \textbf{经典渐近推断与统计假设检验方法的系统应用}：应用经典大样本渐近理论~\cite{lehmann1998theory,vandervaart1998asymptotic}、一阶 Delta 方法与非参数 Bootstrap~\cite{efron1994introduction}，计算温度估计量 $\hat T$ 的渐近方差与三类置信区间（Wald 区间、单调逆变换反演区间及 Bootstrap 百分位区间）；应用经典的 Wald、Score 与 LRT 检验对 $H_0: T=1$ 进行统计显著性裁决，并结合 Platt~\cite{platt1999probabilistic} 提出的双参数 Platt Scaling 对照，利用 McNemar 检验实证比较两类方法的准确率保序性。
\item \textbf{基于无泄漏协议的实证对比与跨模型深度效应分析}：在 CIFAR-10 官方测试集上采用分层无泄漏切分方案（$2{,}000:8{,}000$），完成 ResNet-20 与 ResNet-56 的独立实验与统计检验；在 $p < 10^{-16}$ 的显著性水平下证实了两类残差网络的过度自信现象，揭示了深层网络（ResNet-56, $\hat T=1.6904$）相比浅层网络（ResNet-20, $\hat T=1.6245$）更为显著的校准需求，并在独立评估集上实现了 ECE 下降超过 $66\%$ 的校准改善。
\item \textbf{有限样本实证现象的统计机理解释}：利用已有统计理论对实验中观察到的现象作出解释：用原假设点与 MLE 估计点处的局部 Fisher 信息差异（$I_n(1) \ll I_n(\hat\beta)$）解释三大检验统计量 $\text{Score} > \text{Wald} > \text{LRT}$ 的经验排序；借助抽样理论中的有限总体修正（FPC）公式解释样本量扩增下经验标准差略低于超总体 Delta 理论标准误的成因；并通过双对数线性回归验证了估计量满足经典 $O(n^{-1/2})$ 渐近相合速率。
\end{enumerate}

\subsection{论文结构安排}

全文共分为四个主体部分：
第~\ref{sec:intro} 节界定概率校准的科学问题，梳理国内外学术文献脉络并确立本文研究目标；
第~\ref{sec:theory} 节建立 Temperature Scaling 与 Platt Scaling 的统计模型设定，推导极大似然估计、Score 方程与二阶 Fisher 信息量，证明对数似然严格凹性，并构建 Delta 渐近推断与假设检验体系；
第~\ref{sec:empirical} 节详细介绍无泄漏实验设计、预处理与描述统计，围绕图~\ref{fig:fig1_eda} 至图~\ref{fig:fig5_convergence} 展开探索性数据分析、参数推断、综合评估及渐近理论的实证检验；
第~\ref{sec:discussion} 节深入剖析深度效应传导机制、三大检验排序成因、有限总体修正效应以及单双参数保序权衡，指出研究局限并给出标准化评测建议；
文末附有规范的 GB/T 7714 参考文献著录及附录。
